\documentclass[12pt]{article}
\usepackage[utf8]{utf8}
\usepackage[margin=2cm]{geometry} % Margens ligeiramente reduzidas para garantir 1 página
\usepackage{booktabs}
\usepackage{tabularx}
\usepackage{array}

\begin{document}

% --- TÍTULOS CENTRALIZADOS ---
\begin{center}
    % Nome da Disciplina em 16pt (bold)
    {\fontsize{16pt}{20pt}\selectfont \textbf{Teoria do Aprendizado Estatístico / PI 3}}\\[0.3cm]
    % Nome do Professor em 12pt (bold)
    {\fontsize{12pt}{15pt}\selectfont \textbf{Professor:} João Paulo Ferreira de Mello}\\[0.2cm]
    % Nome dos Integrantes em 12pt
    {\fontsize{12pt}{15pt}\selectfont \textbf{Integrantes:} Ivan Guimarães, Jorge Sugano, Gustavo}\\[0.5cm]
\end{center}

% --- DESCRIÇÃO DO DATASET (PARÁGRAFO JUSTIFICADO) ---
\noindent
\textbf{Descrição do Dataset:} O presente estudo utiliza a base de dados do \textbf{Índice Paulista de Desenvolvimento Municipal (IPDM)}, elaborada pela Fundação Seade. Trata-se de um indicador sintético desenvolvido para avaliar e acompanhar o desempenho socioeconômico dos 645 municípios do Estado de São Paulo em periodicidade bianual (no período de 2014 a 2024). A base mensura as condições municipais por meio de um índice geral e de três dimensões fundamentais: Riqueza, Longevidade e Escolaridade, fornecendo insumos para a análise comparativa da evolução do desenvolvimento humano e econômico no estado.

\vspace{0.4cm}

% --- TABELA DO DICIONÁRIO DE DADOS ---
\begin{table}[htbp]
\centering
\fontsize{10pt}{13pt}\selectfont % Fonte 10pt com entrelinha ajustada
\renewcommand{\arraystretch}{1.15} % Espaçamento ideal para caber na página

\caption{Dicionário de Dados --- Índices e Indicadores do IPDM (2014--2024)}
\label{tab:dicionario_ipdm}
\vspace{0.2cm}

\begin{tabularx}{\textwidth}{>{\ttfamily}l l l X X}
\toprule
\rule{0pt}{2.5ex}\textbf{Variável} & \textbf{Tipo} & \textbf{Domínio / Variação} & \textbf{Descrição} & \textbf{Fonte} \\[0.3ex]
\midrule
cod\_ibge    & Numérico & [3500105, 3557303]    & Código oficial de 7 dígitos do IBGE para o município & IBGE \\
municipio   & Texto    & Nomes de municípios   & Nome do município do Estado de São Paulo            & IBGE \\
valor       & Numérico & [0,000; 1,000]        & Nota/pontuação da cidade no IPDM geral ou na dimensão & Seade \\
ano         & Numérico & [2014, 2024]          & Ano de referência dos dados (série bianual)          & Seade \\
tipo        & Texto    & Categórico            & Escopo da linha (Geral, Riqueza, Longevidade ou Escolaridade) & Seade \\
valor\_estado & Numérico & [0,000; 1,000]        & Pontuação equivalente para o total do Estado de SP  & Seade \\
indicador1  & Numérico & Variável por dimensão & 1º componente específico da dimensão de análise      & Seade \\
indicador2  & Numérico & Variável por dimensão & 2º componente específico da dimensão de análise      & Seade \\
indicador3  & Numérico & Variável por dimensão & 3º componente específico da dimensão de análise      & Seade \\
indicador4  & Numérico & Variável por dimensão & 4º componente específico da dimensão de análise      & Seade \\
indicador5  & Texto / Num. & Variável por dimensão & Identificador ou nota do indicador da dimensão       & Seade \\
\bottomrule
\end{tabularx}
\end{table}

\end{document}